\textbf{Data generation.} Conformations of $N=16$ residues are enumerated by depth-first search with the first step fixed to $+x$ and the first turn fixed to $+y$ (removing the symmetry group of the square lattice), giving \rhval{data/enumeration/hp16/n_conformations/mean} conformations. All $2^{16}$ sequences are scored exactly. Of these, \rhval{data/enumeration/hp16/n_unique_gs_sequences/mean} have a unique ground state, and those ground states cover \rhval{data/enumeration/hp16/n_designable_conformations/mean} conformations; these are the \emph{designable} conformations and form the set of possible targets (a target that no sequence designs would make the task unsolvable by construction). Distinct designable conformations are the unit of the split.

\textbf{Splits.} Reading a chain backwards maps a conformation and its uniquely folding sequence to another conformation and its reversed sequence. If a conformation were in the training set and its reverse in the test set, the test design would be a trivial transformation of a training pair; an early version of our split that ignored this leaked and gave inflated test success, so all splits are by chain-reversal class (a conformation and its reverse are kept together). A fixed, seed-independent 20\% of classes is the DEV set, used only for tuning. For each of five seeds (0--4) the remaining classes are shuffled with that seed, 25\% become test targets (about 90 conformations) and 75\% are training conformations (about 276). Within a seed all systems see identical test targets. The training pairs are those whose sequence lies in a randomly sampled fraction of the $2^{16}$ sequence space (fraction 1.0 unless stated; a sweep varies it) and whose ground state is a training conformation; this makes the data ``sampled sequences scored by the oracle''.

\textbf{Budgets and metrics.} $K\in\{1,10,100,1000\}$. Primary metric: succ\_k$K$ (Section~\ref{sec:method}); per-regime variants split the targets into those with exactly one uniquely folding sequence (``single'') and those with several (``multi''). We also log the energy gap of the returned design (floor of $f$), the held-out per-residue NLL of the true designing sequences under the AR model, and the fraction of all sampled designs that succeed.

\textbf{Hyperparameters and tuning.} AR designer: width 64, 4 heads, feed-forward width 256, dropout 0.1, two encoder and two decoder layers, AdamW (learning rate $2\times10^{-3}$, one-cycle schedule, weight decay 0.01), batch 64, 60 epochs, $\tau=0.7$. Epochs $\in\{30,60\}$ and $\tau\in\{0.7,1.0\}$ (four configurations) were chosen on DEV (seed 100) by K=1 success; the other settings were fixed in advance. Simulated annealing: $T_0\in\{0.5,1,2\}\times T_1\in\{0.05,0.2,0.5\}$ (nine configurations) on DEV, chosen by the sum of success at $K=100$ and $K=1000$: $T_0=0.5$, $T_1=0.2$, reused unchanged for the heuristic-initialised and graded-objective variants except that the graded-objective variant had its own nine-configuration DEV grid (same choice). After the audit, six further low-temperature SA configurations per objective ($T_0\in\{0.1,0.25,0.35\}$) were run on DEV; none improved the selection score, so $T_0=0.5$, $T_1=0.2$ was kept (see Limitations). AR and SA thus received comparable, small tuning budgets. Ablation variants reuse the tuned configuration without re-tuning.

\textbf{Hardware and runtime.} A shared multi-core CPU machine, two threads per process, no GPU. An AR run (training plus sampling for all test targets) takes a few minutes; annealing and random search take seconds. Oracle calls are table look-ups, so ``budget'' counts evaluations of the objective, not wall-clock. The code is in \texttt{method/}, the run driver in \texttt{experiments/run\_all.sh}.

\textbf{Statistics.} Five seeds; we report mean $\pm$ standard deviation over seeds and paired two-sided $t$-tests across seeds (same seed means the same test targets). With five seeds these tests have low power and no multiplicity correction is applied.
